% Chapter 12: Tesis de Inversión y Proyecciones Financieras de REDCO
\chapter{Tesis de Inversión y Proyecciones Financieras}

\subsection{Análisis de Escenarios Financieros}

Para modelar la viabilidad de la estrategia, se proponen tres escenarios financieros a 5 años para REDCO, considerando la exclusión del mercado ruso y la penetración en Chile, Perú y USA:

\begin{enumerate}
    \item \textbf{Escenario A: Base Ajustado (Exclusión de Rusia, Foco Chile Core):}
    \begin{itemize}
        \item Crecimiento anual de ingresos: 5.0\%.
        \item Margen EBITDA: 22\%.
        \item Conversión de caja (Caja Real / Devengo): 92\%.
        \item Dependencia del fundador: Disminuye a 15\% al año 5.
    \end{itemize}

    \item \textbf{Escenario B: Crecimiento Acelerado (Penetración Perú + USA + Alianza Tecnológica):}
    \begin{itemize}
        \item Crecimiento anual de ingresos: 12.5\%.
        \item Margen EBITDA: 28\% (por prima tecnológica).
        \item Conversión de caja: 95\%.
        \item Dependencia del fundador: Disminuye a $<5\%$ al año 5.
    \end{itemize}

    \item \textbf{Escenario C: Conservador / Downside (Estancamiento comercial, no expansión):}
    \begin{itemize}
        \item Crecimiento anual de ingresos: 1.5\%.
        \item Margen EBITDA: 15\% (compresión por competencia de precios).
        \item Conversión de caja: 85\% (retrasos de pago en Chile/Perú).
        \item Dependencia del fundador: Se mantiene alta ($>40\%$).
    \end{itemize}
\end{enumerate}

\begin{table}[htbp]
\centering
\caption{Proyecciones Financieras Comparativas por Escenario (USD Millones)}
\begin{tabular}{@{}lrrrrrr@{}}
\toprule
\textbf{Escenario / Métrica} & \textbf{2026} & \textbf{2027} & \textbf{2028} & \textbf{2029} & \textbf{2030} & \textbf{2031} \\
\midrule
\textbf{Escenario A (Base)} & & & & & & \\
Ingresos & \$3.20 & \$3.36 & \$3.53 & \$3.70 & \$3.89 & \$4.08 \\
\rowcolor{tablealt} EBITDA (22\%) & \$0.70 & \$0.74 & \$0.78 & \$0.81 & \$0.86 & \$0.90 \\
Caja Disponible & \$0.64 & \$0.68 & \$0.72 & \$0.75 & \$0.79 & \$0.83 \\
\midrule
\textbf{Escenario B (Acelerado)} & & & & & & \\
Ingresos & \$3.20 & \$3.60 & \$4.05 & \$4.56 & \$5.13 & \$5.77 \\
\rowcolor{tablealt} EBITDA (28\%) & \$0.90 & \$1.01 & \$1.13 & \$1.28 & \$1.44 & \$1.62 \\
Caja Disponible & \$0.85 & \$0.96 & \$1.07 & \$1.21 & \$1.36 & \$1.53 \\
\midrule
\textbf{Escenario C (Downside)} & & & & & & \\
Ingresos & \$3.20 & \$3.25 & \$3.30 & \$3.35 & \$3.40 & \$3.45 \\
\rowcolor{tablealt} EBITDA (15\%) & \$0.48 & \$0.49 & \$0.50 & \$0.50 & \$0.51 & \$0.52 \\
Caja Disponible & \$0.41 & \$0.42 & \$0.42 & \$0.43 & \$0.43 & \$0.44 \\
\bottomrule
\end{tabular}
\label{tab:financial_projections}
\end{table}

\subsection{Conclusión y Tesis de Valor}

La tesis de inversión para el crecimiento de REDCO se sustenta en que la diversificación geográfica controlada y la transformación del modelo operativo ( Execution Premium) neutralizan el riesgo de concentración geopolítica (Efecto Seligdar/Rusia) y aumentan sustancialmente el valor del capital de la firma. Al desligar la ejecución técnica de la persona del fundador, REDCO pasa de ser una práctica consultora personal de alto riesgo a convertirse en una empresa institucional escalable, atractiva para potenciales fusiones o adquisiciones en el horizonte de 2031.
